% Chapter: The Stories
% Track: 2
% Content: Kangaroos (from HealerStories archive)

\settrack{tracktwo}

\chapter{The Stories}
\label{ch:t2-stories}
\label{pos05:the-stories}

% TODO: epigraph

\begin{quote}\small\textit{%
``We cultivated our land, but in a way different from the white man. We endeavored to live with the land; they seemed to live off it. I was taught to preserve, never to destroy.'' --- Tom Dystra, Australian Aboriginal Elder
}\end{quote}

\vspace{0.5cm}

The year is 1978.  The place is a large quarter horse ranch in New South Wales, Australia.

A twelve year old boy rides the fence line astride a four year old quarter horse mare.  It's a three day ride around the ranch.  A working ranch never has enough hands.  Securing the fence line is a man's job.  This boy has become a young man.  David rides the fence line with water, food, a knife, a first aid kit, basic field gear, a fence repair kit, and a .303 Lee Enfield rifle.  He knows how to operate a HAM radio but today he has no radio. Any problems along the fence line must be fixed or reported.

Along the way the young man will visit various people who live on the ranch.  Some are retired Australian military men and their families.  Most of the people he will visit are of Kuringai aboriginal ethnicity. Many Kuringai people live on this land as they have for thousands of years.  They maintain good relations with the ranch `owners', who protect them from other white Australians and allow them to live there in peace. In return the Kuringai perform certain maintenance and security responsibilities. These responsibilities include training the children of the ranch owner in aboriginal bushcraft.

David senses trouble before he rides up to the North Gate.  Subtle signs suggest something is amiss.  His friends Koori and Ngemba, both aboriginal adult men, show themselves as he approaches on horseback.  He rides up to them and dismounts.  Friendly greetings are brief. Trouble is in the air.

Ngemba and Koori confirm the boy's suspicions. Ngemba tells David that a ute, or pickup truck, loaded with gubbins, dangerous armed white men, have crashed the gate and are on the ranch making trouble.  They are drinking beer and shooting kangaroos with their rifles. Kangaroos are sacred to the Kuringai.  Killing kangaroos this way is illegal but rarely punished. Ngemba and Koori have already spread the alarm.

These gubbins are known.  They have trespassed on the ranch before.  Each time they leave behind empty bottles of liquor, bullet casings, and dozens of dead 'roos.  The Kuringai people weep at the waste but can do nothing.  The unruly white men have rifles so the Kuringai must avoid a direct confrontation.  None of the Kuringai will show themselves to the gubbins.   Whenever aborigines fight whites always the abos are always punished or killed, regardless of who started it, who won, or who has the moral high ground or law on their side.  The white men in the ute will be long gone before the police can arrive.

As the son of the ranch owner David has much higher status with the police than do the aborigines.  If anyone is to do something it's up to him.

Koori tells David that the gubbins are parked about five kilometers away on the north side of Wattle Ridge. Ngemba's friend Darawal, who can effect a very convincing white man's accent, is already on his way to the nearest telephone to alert the police.  This land is big, though, and it will take police at least three hours to arrive.  None of the retired SASR veterans who live on the ranch happen to be nearby.

The police are on their way.   While David is not Kuringai, he and his twin sister grew up around them and learned their ways. Both he and his sister lived among them for months at a time.  Aborinal hunters taught them stealth, tracking, and bushcraft.  David is half Maori on his mother's side \& half European Australian on his father's side.  Most importantly, his parents own the ranch.  He has social and legal authority to defend the ranch against invaders in a way the Kuringai do not.  He can get away with things that would get aborigines arrested or killed.

With Koori and Ngemba's help David forms a daring plan. He will try to delay the gubbins until the police arrive.  Ngemba and Koori wish him luck.  They will stealthily observe from a safe distance.

David checks his gear and saddles up.  He stays off the road and circles towards the intruders, careful to keep hills between them.  Sporadic gunfire tells him where they are.  David has hunted before but this is his first time hunting men.  Their rifles make them very dangerous.  Their probable drunkenness makes them no less dangerous.

The boy knows this hilly land.  He walked it last year with Koori.  He rides up a rocky hill to a ridgeline.  The trespassers are on the other side.  He dismounts about 600 meters from the dangerous armed intruders, on the other side of a hill, and leads his horse. Pop! Pop! More gunfire as another kangaroo dies.   He feeds and waters his horse then leaves it hobbled and tied. Horse and boy are ready for a quick departure along a covered escape route.

David cautiously approaches the ridgeline.   The youth suppresses his fear and controls his breathing. He readies his binoculars and checks his rifle.  He crawls the last few meters to the ridgeline and peeks between big rocks towards the intruders below.

The dry valley far below is a scene of violence.  The ute is a big four seater, what most non- Australians would call a large four-by-four pickup truck.   Four men equipped with rifles and beer are shooting kangaroos.  About 20 'roos lie dead or dying.  Red blood still flows.  David watches through binoculars.  One man raises his rifle, shoots another 'roo, then swigs more beer.  David feels his blood rise.  He lets loose a quiet Kuringai curse at these awful men. He will not kill them but he is determined to delay them until the police arrive.

The boy has a clear line of fire down on the intruders and their vehicle.  They are about 400 meters away and still unaware of his presence.  David again checks his rifle and spare ammo clips.   He carries a powerful bolt action .303 Lee Enfield rifle. The Lee Enfield was the primary battle rifle of the British Army through two world wars. His uncle, a special forces SAS combat veteran who fought in several wars, taught him to shoot this rifle. A powerful .303 bullet, besides being lethal to most living things, is an effective anti-material round.  A .303 bullet will destroy an engine block, flatten tires, and can otherwise disable a vehicle. David takes careful aim at the big ute far below.  400 meters is a long range shot through iron sights but the engine is a big target.  He will have to shoot past a rifleman but judges he can safely do so without hitting the man.  BANG! goes his first round.  The boy palms the bolt, reloads, and keeps firing. It takes the intruders a little while to realize someone is shooting at them and more time to recover from their initial shock and panic.  David's mad minute of shooting puts 5 bullets through the ute's engine block and flattens multiple tires.  That ute isn't driving away anytime soon.

The intruders now have a disabled vehicle, are in unfriendly territory, don't know the terrain, and lack adequate supplies to hike out.  They still manage to locate the shooter and return fire.  David feels bullets crack the air above his head and hears more bullets strike the rocks nearby. A small chip of rock lands harmlessly on his leg.  David's blood is up and he can clearly see one of the men shooting at him. David is tempted to kill the man but decides against such murderous foolishness. Instead he stops shooting. Still prone, he crawls back out of sight.  They continue to fire sporadically at his previous location.  David withdraws, recovers his horse, then rides away and circles around the intruders.

Twnety minutes later, from a distant ridgeline, the boy cautiously and stealthily observes the angry armed intruders. They do not spot him. The intruders are not prepared to escape on foot and their car's engine is scrap metal. David does not reveal his location and fires no more bullets that day.  He patiently observes.

By the time police arrive the intruders are glad to be rescued, even by police.   The police check everyone for safety, provide the intruders with water, then arrest them.  David keeps to his hidden vantage point and watches.   The intruders are handcuffed and locked in police vehicles.  Just before they drive away one police officer gets out of his car, looks towards the ridgeline where David watches from concealment, gives a gesture of approval, then gets back in his car and drives away with the crims.  David rides away, reports the action to Koori, and continues the serious job of riding the fence line of his parents' ranch.

A few days later the police visit David's parents.  The no-longer-a-boy accurately described the entire incident to his mum and dad.  The version told by police matches his story.  The police laud the lad for helping them catch the criminals, who are still in jail. They'll be released soon but are unlikely to return. The wrecked ute is towed away for scrap.

The same police officer who gestured his approval, we'll call him Officer Kevin, wants to meet the lad.  They meet.  Officer Kevin, who is a military veteran, begins to keep tabs on David.  Kevin become David's friend and mentor.  Years later, when teenage David gets in trouble with the law, Officer Kevin is there to help.

\chapterreturn
